\documentclass[a4paper,12pt]{article}
\usepackage[italian]{babel}
\usepackage[utf8]{inputenc}
\usepackage[T1]{fontenc}
\usepackage{geometry}
\usepackage{setspace}
\usepackage{eso-pic}
\usepackage{xcolor}

\geometry{margin=2.5cm}
\onehalfspacing

% Cambia SOLO questa riga per creare C1, C2, C3
\newcommand{\VersionCode}{C1}

\AddToShipoutPictureBG{
	\AtPageCenter{
		\makebox(0,0){
			\rotatebox{30}{
				\textcolor[gray]{0.85}{
					\fontsize{18}{18}\selectfont
					\begin{tabular}{cccc}
						TPSIT 4AIQ \VersionCode\ 02/03/26 & TPSIT 4AIQ \VersionCode\ 02/03/26 & TPSIT 4AIQ \VersionCode\ 02/03/26 & TPSIT 4AIQ \VersionCode\ 02/03/26 \\
						TPSIT 4AIQ \VersionCode\ 02/03/26 & TPSIT 4AIQ \VersionCode\ 02/03/26 & TPSIT 4AIQ \VersionCode\ 02/03/26 & TPSIT 4AIQ \VersionCode\ 02/03/26 \\
						TPSIT 4AIQ \VersionCode\ 02/03/26 & TPSIT 4AIQ \VersionCode\ 02/03/26 & TPSIT 4AIQ \VersionCode\ 02/03/26 & TPSIT 4AIQ \VersionCode\ 02/03/26 \\
						TPSIT 4AIQ \VersionCode\ 02/03/26 & TPSIT 4AIQ \VersionCode\ 02/03/26 & TPSIT 4AIQ \VersionCode\ 02/03/26 & TPSIT 4AIQ \VersionCode\ 02/03/26 \\
					\end{tabular}
				}
			}
		}
	}
}


\title{Verifica di TPSIT - 4AIQ - Compito C}
\author{Prof. Fedeli Massimo}
\date{02/03/26}

\begin{document}

\maketitle
Cognome e Nome:

\vspace{0.15cm}
\noindent\textbf{Codice versione: \VersionCode}

\section*{Compito}

Hai a disposizione il seguente schema logico; svolgi l'esercizio 1 oppure l'esercizio 2 a tua scelta e rispondi alle domande di teoria
argomentando in modo approfondito:

\begin{verbatim}
	EVENTI(id PK, titolo, luogo, data_evento, ora_inizio, ora_fine, capienza)
	PARTECIPANTI(id PK, cognome, nome, email UNIQUE, telefono)
	BIGLIETTI(id PK, evento_id FK, partecipante_id FK, codice UNIQUE,
	prezzo, stato ENUM('prenotato','pagato','annullato'),
	data_acquisto)
	ACCESSI(id PK, biglietto_id FK, timestamp_accesso, esito ENUM('ok','negato'),
	motivo)
\end{verbatim}

\section*{Esercizio 1}

\subsection*{Obiettivo}

Realizzare una piccola applicazione web che consenta di inserire un nuovo partecipante
nella tabella \texttt{PARTECIPANTI} tramite una pagina web (client) e un endpoint PHP
(backend) che salvi i dati su MySQL. (Host: 127.0.0.1 \ \ DBNAME: Eventi)

\subsection*{Vincoli e requisiti tecnici}

\begin{itemize}
	\item Frontend: una pagina \texttt{index.html} con un form e uno script JavaScript. \textbf{2.5 pt}
	\item Backend: un file PHP (esempio: \texttt{api\_partecipanti\_create.php}) che riceva i dati e inserisca la riga nel database. \textbf{2.5 pt}
\end{itemize}

Requisiti: 
\begin{itemize}
	\item Comunicazione: il client deve inviare i dati al backend usando \texttt{fetch} in POST, in formato JSON.
	 \item Validazione minima: email obbligatoria e in formato valido; nome e cognome obbligatori; telefono facoltativo.
	\item Sicurezza: il backend deve usare query parametrizzate (PDO con prepared statements).
\end{itemize}

\section*{Esercizio 2}

\subsection*{Obiettivo}

Realizzare una piccola applicazione web che consenta di visualizzare l'elenco dei
biglietti emessi per tutti gli eventi, includendo anche le informazioni minime
dell'evento e del partecipante, e lo stato del biglietto. (Host: 127.0.0.1 \ \ DBNAME: Eventi)

\subsection*{Vincoli e requisiti tecnici}

\begin{itemize}
	\item Frontend: una pagina HTML che mostri i risultati in una tabella e uno script JavaScript. \textbf{4 pt}
	\item Backend: un file PHP (esempio: \texttt{api\_biglietti\_list.php}) che legga i dati dal database e restituisca un JSON. \textbf{4 pt}
\end{itemize}

Requisiti:
\begin{itemize}
	\item Comunicazione: il client deve richiedere i dati al backend usando \texttt{fetch} e ricevere dati in formato JSON.
	\item Query: il backend deve eseguire una SELECT con JOIN tra:
	\begin{itemize}
		\item \texttt{BIGLIETTI} e \texttt{PARTECIPANTI}
		\item \texttt{BIGLIETTI} e \texttt{EVENTI}
		\item \texttt{BIGLIETTI} e \texttt{ACCESSI} (LEFT JOIN)
	\end{itemize}
	\item Campi minimi da mostrare in tabella: codice biglietto, titolo evento, data\_evento, luogo, nominativo partecipante, email, stato biglietto, data\_acquisto, ultimo \texttt{timestamp\_accesso} ed \texttt{esito}.
	\item Sicurezza: il backend deve usare PDO e prepared statements anche per eventuali parametri di filtro.
\end{itemize}

\section*{Domande}

\textbf{Domanda 1 – Cookie}
Spiega cosa sono i cookie in ambito web, descrivendo a cosa servono, dove vengono memorizzati e come vengono utilizzati dal protocollo HTTP per mantenere informazioni tra richieste successive. Illustra inoltre la differenza tra cookie di sessione e cookie persistenti, indicando almeno un possibile utilizzo pratico di ciascuno.
\textbf{Domanda 2 – Sessioni}
Descrivi il meccanismo delle sessioni nelle applicazioni web, spiegando perché sono necessarie e come permettono di mantenere lo stato tra richieste HTTP. Illustra il ruolo del session ID, dove vengono memorizzati i dati di sessione e come avviene il collegamento tra client e server. Indica infine le principali differenze rispetto ai cookie in termini di sicurezza e gestione dei dati.

\end{document}